%*******************************************************************
% Chapter 9: Conclusion
%*******************************************************************

We presented UniAdapFuse, a reliability-aware multimodal fusion framework for UAV disaster perception, together with a physics-guided Hazard Prior Gate for smoke and fire detection. We evaluated twenty-six classification configurations on a 72{,}997-row corpus assembled from six public sources, seven detectors on the 21{,}527-image D-Fire benchmark, and a prototype joint graph. Throughout, our analysis distinguishes measured RGB from RGB-derived pseudo-thermal maps, label-matched audio from synchronized capture, and scene labels from bounding-box localization. It further separates within-collection from cross-sequence performance: retraining under a capture-group-aware split scores 9.60 percentage points lower in combined macro-F1, and decomposing those predictions by source shows that source identity and label are entangled in the assembled corpus.

The single most consequential finding is that the near-ceiling scores are inflated by row-level splitting and cannot be separated from source-label confounding. Retraining AdapFuse-v1 under a capture-group-aware split drops combined macro-F1 from 98.74\% to 89.14\% ($-9.60$ pp, single completed seed), with disaster classification losing 14.28 pp against the victim task's 4.93 pp. That retraining also disabled class balancing and used batch size 2, so the drop mixes leakage with recipe effects; a controlled rerun with the row-level recipe is needed to isolate the leakage component. The per-source decomposition (Section~\ref{sec:per_source}) shows that within-source accuracy is reasonable (0.82--0.998 for disaster), and that the low per-source macro-F1 values written by the evaluation script arise because classes absent from a source are scored as zero. It also shows the underlying problem: victim presence is fully determined by the source dataset in the grouped test set, SARD contains only clean scenes and FireNet only fire, so a model that recognises the dataset can score well without recognising people or hazards. The current experiments cannot rule this out, which qualifies every row-level number summarised below.

Within that qualification, the visual backbone variants are strong scene classifiers on the present split, with EfficientNet-B0 reaching 99.24\% combined F1 and MobileViT-XXS reaching 98.99\%; the victim metric is binary presence classification, and no person box, object size, or localization accuracy is measured. Fusion itself does not raise row-level F1: the RGB-only baseline (98.85\% combined F1) is 0.11 points above AdapFuse-v1 (98.74\%), and the no-RUE variant (98.75\%) matches it, so the reliability-aware design is supported here as an architecture that fuses heterogeneous inputs and runs with missing modalities, not as a source of accuracy gains. Retraining each single-modality baseline with three backbones confirms this: RGB-only MobileViT-XXS (98.99\%) equals the fused model with the same backbone, and with the backbone held fixed fusion adds at most 0.36 points. Trained through the same pipeline under published protocols, RGB-only EfficientNet-B0 exceeds the reported results on AIDER (94.00\% macro-F1 against 88.5\%) and on FLAME's cross-camera test (83.72\% accuracy against 76.23\%). On FLAME, however, every model scores about 99.9\% on validation frames from the training drone and 17--37 points lower on the test drone, the same gap between within-capture and new-capture performance that the grouped split exposes. Among the compared fusion runs, AdapFuse-v1 attains the lowest loss, with Design E reporting loss 0.0020 and ECE 0.0337, although its 0.73\% false-positive rate is computed on the row-level scene-classification split and cannot be attributed to the nuisance head, whose loss is zero without nuisance labels. The component variants identify hypotheses for controlled follow-up rather than settled ablation results: relative to separate no-GRU, no-attention, and no-quality-feature runs, the base result is higher by 0.60, 0.35, and 0.34 combined-F1 points, but the repository has no repeated seeds or confidence intervals, the GRU receives sequence length one, and the nuisance comparison cannot demonstrate nuisance regularization because no nuisance targets are supplied.

The label-matched audio channel is not merely weak but leaky. The metadata reuse 2{,}000 ESC-50 clips across 39{,}003 assignments without synchronized UAV capture, and on every row with audio the clip's ESC-50 class identifies fire/smoke versus clean. The audio-only baseline's low 38.2\% disaster macro-F1 comes from missing audio on 47\% of test rows and from the absence of audio for the collapse and other-hazard classes, not from uninformative clips. The row-level model's audio gate settles near zero ($\approx$0.011 on rows with audio), whereas the grouped retraining opens it fully (1.0), consistent with that run using the shortcut. Data alignment, not acoustic backbone choice, is therefore the main limitation; because the intended AudioSet-pretrained encoder loaded no convolution weights, pretrained audio transfer was not tested either. The two-channel metadata variants remain competitive, with RGB plus pseudo-thermal reaching 98.50\% combined F1 and pseudo-thermal plus audio reaching 98.05\%, but because pseudo-thermal is derived from RGB these experiments do not validate zero-visibility or independent thermal-sensor operation.

On the joint and dense tasks, the revised joint recipe recovers classification in the prototype but not detection: UniAdapFuse v2 reaches 98.57\% combined F1 and 99.26\% disaster accuracy after the v1 classification collapse, yet its saved evaluation reports detection $\mathrm{mAP}_{50}=0.0$. The recovery combines stop-gradient isolation with distillation from AdapFuse-v1 and several other recipe changes, so it cannot be credited to isolation alone, and both joint runs show widespread non-finite losses, so dense joint detection remains unvalidated. Physical prior gating did not help on D-Fire: HPG YOLO26n reaches 70.53\% $\mathrm{mAP}_{50}$ in 8.65 hours against 71.47\% in 4.45 hours for baseline YOLO26n under the same 40-epoch recipe, with near-identical validation curves. Knowledge distillation, finally, yields only modest compression: AdapFuse-v2 halves the teacher's projection dimension with negligible accuracy loss ($<$0.2 percentage points combined F1), but because the dominant ResNet-18 and MobileNetV3-Small backbones are untouched, the measured total parameter reduction is 12.94M $\to$ 12.30M (4.9\%) rather than proportional to the projection-width halving. The teacher AdapFuse-v1 measures 14.32 ms per frame (69.9 FPS) on the evaluation workstation GPU, and dedicated benchmarking of AdapFuse-v2 on physical embedded edge hardware is identified as future work.

Reconciling these outcomes with the requirements set out in Chapter~3 makes several revisions explicit rather than leaving them as a silent gap between chapters, as is normal in an iterative research project. Chapter~3 called for at least 6{,}000 \textit{synchronised} (RGB, thermal, audio) sequences; our dataset instead pairs ESC-50 audio to visual samples by disaster label rather than by simultaneous capture, and every populated thermal path points to an RGB-derived pseudo-thermal map rather than a sensor measurement, so the near-constant audio gate ($r_{\text{audio}} \approx 0.006$) shows that the fitted model largely ignores this channel rather than proving calibrated reliability estimation. Chapter~3 also specified at least 95\% \textit{detection} accuracy for wildfire scenes and at least 90\% for USAR victim \textit{detection}; against these targets our system reports two distinct tasks, since scene-level \textit{classification} (99.4\% disaster accuracy, 98.36\% victim F1, Tables~\ref{tab:adaptfuse_extended} and~\ref{tab:main_results}) exceeds both thresholds on the row-level split while dense bounding-box \textit{detection} on D-Fire (Table~\ref{tab:dfire_results}) reaches 77.10\% mAP$_{50}$ for smoke/fire and no person/survivor bounding-box benchmark is reported (Section~\ref{sec:object_detection}, scope note). The thresholds were therefore met for row-level classification but not for dense object detection, and the two metrics should not be conflated. Relatedly, the architecture's Survivor RoI detection head (Section~\ref{sec:uniadaptfuse_arch}) is structurally implemented but never evaluated against ground-truth person bounding boxes in this work; only victim \textit{presence} classification is benchmarked, leaving person-level localisation evaluation, for example on SARD, to future work. Table~\ref{tab:req_trace} summarises the reconciliation.

\begin{table}[H]
\centering
\caption{Chapter~3 requirements against the evidence reported in Chapter~5.}
\label{tab:req_trace}
\resizebox{\textwidth}{!}{%
\begin{tabular}{@{}>{\raggedright\arraybackslash}p{3.0cm}>{\raggedright\arraybackslash}p{3.0cm}>{\raggedright\arraybackslash}p{4.4cm}>{\raggedright\arraybackslash}p{2.4cm}>{\raggedright\arraybackslash}p{4.6cm}@{}}
\toprule
\textbf{Requirement} & \textbf{Specification} & \textbf{Reported evidence} & \textbf{Status} & \textbf{Qualification} \\
\midrule
Wildfire detection & $\geq$95\% detection accuracy & 99.44\% disaster-classification accuracy (Table~\ref{tab:adaptfuse_extended}); 77.10\% smoke/fire mAP$_{50}$ on D-Fire (Table~\ref{tab:dfire_results}) & Met for classification only & Scene classification on the row-level split; grouped split gives 89.05\% accuracy; bounding-box detection is a different metric \\
\addlinespace
USAR victim detection & $\geq$90\% detection accuracy & 98.36\% victim-presence macro-F1 (Table~\ref{tab:adaptfuse_extended}) & Met for classification only & Presence label is determined by source dataset; no person bounding boxes evaluated \\
\addlinespace
False-positive rate & $\leq$8\% & 0.73\% row-level scene FPR (Table~\ref{tab:adaptfuse_extended}) & Met on row-level split & Not an operational field false-alarm rate; nuisance head unsupervised \\
\addlinespace
Graceful degradation & $\geq$85\% accuracy with any one modality missing & Two-sensor subsets 98.05--98.50\% combined F1 (Section~\ref{sec:two_sensor}) & Partially met & Pseudo-thermal derives from RGB, so losing RGB is not an independent-sensor test \\
\addlinespace
Latency & $\leq$500\,ms at $\leq$15\,W & 14.32\,ms on RTX 5090 workstation (Table~\ref{tab:latency_profile}) & Not validated & No embedded-hardware latency or power measured \\
\addlinespace
INT8 quantisation & All encoders & Not implemented & Not met & Future work \\
\addlinespace
Synchronised dataset & $\geq$6{,}000 synchronised tri-modal sequences & Pseudo-thermal from RGB; label-matched ESC-50 audio & Not met & No synchronised sensor data used \\
\bottomrule
\end{tabular}%
}
\end{table}

Read against that boundary, this work makes four contributions to autonomous UAV aerial perception. The first is the UniAdapFuse prototype, a joint computation graph that combines multi-hazard scene classification, victim-presence classification, and a dense detection branch, in which a revised recipe (phased training, stop-gradient isolation and distillation) recovers classification to 98.57\% combined F1 while the recorded detection $\mathrm{mAP}_{50}$ remains 0.0; the contribution is thus an implemented prototype and a documented failure mode rather than a validated unified detector. The second is the Hazard Prior Gate, a physics-guided module that derives red-excess chromaticity, smoke achromaticity, and edge cues from the input and applies a bounded residual correction before the detector, together with the negative result that it neither raises accuracy (70.53\% against 71.47\% $\mathrm{mAP}_{50}$) nor shortens training (8.65 against 4.45 hours) for YOLO26n on D-Fire. The third is the audit of the reliability and nuisance mechanisms: the tri-modal architecture implements sigmoid modality gates and a nuisance branch, but in our models the gates saturate at 0 or 1, do not move across corruption severity, and differ between the row-level and grouped retraining (audio off in one, on in the other), and no degradation flags or nuisance labels are supplied by the trainer, so the audit establishes the implementation boundary and the experiments needed before dynamic reliability or nuisance suppression can be claimed. The fourth is reproducible comparative evidence, since the repository preserves results for multiple detector and classification configurations, component variants, provenance counts, and workstation latency, with AdapFuse-v2 containing 12.30M parameters against 12.94M for its teacher and the teacher measuring 14.32 ms on an RTX 5090 workstation (32\,GB VRAM), while embedded deployment remains future work.

Six directions follow directly from these limits. The highest priority is real UAV audio and multimodal synchronization: replacing label-matched ESC-50 audio with hardware-synchronized, UAV-embedded acoustic datasets such as DroneAudioset or DREGON is what would fully unlock acoustic complementary signals under rotor noise. Closely related is dynamic pre-backbone reliability estimation, upgrading the RUE to read pre-backbone physical quality signals such as optical flow variance, dark channel prior, or raw SNR, so that intra-modality degradation can be tracked dynamically under heavy smoke and glare. Third, out-of-distribution cross-dataset generalization should be quantified by evaluating the trained multi-modal and object detection models zero-shot on external aerial datasets such as HIT-UAV (aerial thermal imagery) and FLAME~2 (paired RGB and infrared wildfire frames), measuring domain transferability under varying altitudes and geographic biomes. Fourth, onboard flight testing and edge hardware profiling should deploy and profile the distilled AdapFuse-v2 and UniAdapFuse v2 models directly on embedded low-power accelerators aboard an autonomous UAV in real-world disaster simulation exercises. Fifth, a person/survivor bounding-box localization benchmark should evaluate the architecturally-implemented Survivor RoI head against ground-truth person boxes, for example on SARD, to complete the dense-detection evaluation alongside the existing smoke/fire results on D-Fire. Two further gaps follow from the component audit: the zero $\mathrm{mAP}_{50}$ of the joint UniAdapFuse v2 detector must be diagnosed (detection loss weighting, head learning rate, and the target and evaluation path are the first candidates) before joint perception can be claimed, and the GRU must be run on multi-frame sequences with carried-over hidden state before any temporal-modelling claim can be tested.

The sixth direction, source-decorrelated training and evaluation, is the one on which the credibility of every aggregate number depends. Capture-group-aware splitting has now been built and evaluated (Section~\ref{sec:grouped_split}), so near-duplicate leakage is measured rather than merely suspected, but three steps remain. The grouped run must first be repeated with the row-level recipe (class-balanced training, batch size 16), because our run changed both along with the split. The multi-seed interval must then be completed: of three launched seeds one finished cleanly, one selected its checkpoint during a window of degenerate validation, and one halted at epoch 12, so the grouped figure is currently a single-seed point estimate. More importantly, the per-source decomposition in Section~\ref{sec:per_source} shows that source identity and label are confounded in the assembled corpus (most starkly, victim presence is fully determined by source), so sources containing both values of each label, source-balanced class sampling, or leave-one-source-out evaluation are required before any aggregate score can be read as hazard or person discrimination rather than dataset recognition. The audio assignment must likewise be rebuilt so that the clip class no longer encodes the label, for example by pairing each visual row with a clip drawn independently of its label or, better, by using synchronized recordings.
